\chapter{Reproduction and Development}\label{ch:g12:reproduction-development}

\omterm{def:g12:reproduction-development:sexual}{Sexual reproduction} mixes \omterm{prop:g11:dna-replication:genome}{genomes}; \omterm{def:g10:scientific-biology:properties}{development} turns a single \omterm{def:g12:reproduction-development:sexual}{zygote} into
a body with tissues, organs and axes. This chapter follows animal \omterm{def:g12:reproduction-development:sexual}{gametes}
from their production through \omterm{def:g12:reproduction-development:fertilisation}{fertilisation}, sketches the early embryonic
stages of \omterm{def:g12:reproduction-development:cleavage}{cleavage} and \omterm{def:g12:reproduction-development:gastrulation}{gastrulation}, and outlines how placental mammals
nourish an embryo inside the mother.

\section{Sexual reproduction and gametes}

\begin{definition}[Sexual reproduction]
\label{def:g12:reproduction-development:sexual}
\emph{Sexual reproduction}\index{sexual reproduction} involves \omterm{def:g11:meiosis-life-cycles:meiosis}{meiosis}
and the fusion of two \omterm{def:g11:meiosis-life-cycles:ploidy}{haploid} \emph{gametes}\index{gamete} to form a
\omterm{def:g11:meiosis-life-cycles:ploidy}{diploid} \emph{zygote}\index{zygote}. Offspring are genetically distinct
from either parent (barring special cases).
\end{definition}

\begin{definition}[Gametogenesis]
\label{def:g12:reproduction-development:gametogenesis}
\emph{Gametogenesis}\index{gametogenesis} is the formation of \omterm{def:g12:reproduction-development:sexual}{gametes}.
In animals, \emph{spermatogenesis}\index{spermatogenesis} in the testes
produces spermatozoa; \emph{oogenesis}\index{oogenesis} in the ovaries
produces ova (eggs), often with unequal \omterm{def:g11:cell-cycle-mitosis:cytokinesis}{cytokinesis} that conserves
cytoplasm in one large egg.
\end{definition}

\begin{proposition}[Meiosis halves chromosome number]
\label{prop:g12:reproduction-development:meiosis}
\omterm{def:g11:meiosis-life-cycles:meiosis}{Meiosis} produces \omterm{def:g11:meiosis-life-cycles:ploidy}{haploid} nuclei so that \omterm{def:g12:reproduction-development:fertilisation}{fertilisation} restores diploidy.
\omterm{prop:g11:meiosis-life-cycles:prophaseI}{Crossing over} and independent assortment generate \omterm{def:g10:biodiversity-classification:biodiversity}{genetic diversity} among
\omterm{def:g12:reproduction-development:sexual}{gametes}.
\end{proposition}
\admitted

\begin{omfigure}
\begin{tikzpicture}[font=\small]
  \draw[omDef, thick, fill=omDef!15] (0,0) circle (0.55);
  \node at (0,-0.9) {spermatogonium};
  \draw[-Stealth] (0.7,0) -- (1.3,0);
  \draw[omDef, thick, fill=omDef!15] (2.0,0.35) circle (0.4);
  \draw[omDef, thick, fill=omDef!15] (2.0,-0.35) circle (0.4);
  \node at (2.0,-1.0) {meiosis};
  \draw[-Stealth] (2.6,0) -- (3.2,0);
  \foreach \y in {0.55,0.18,-0.18,-0.55}
    \draw[omProp, thick, fill=omProp!20] (3.9,\y) ellipse (0.35 and 0.12);
  \node at (3.9,-1.0) {sperm};
  \begin{scope}[xshift=6.2cm]
    \draw[omMeth, thick, fill=omMeth!15] (0,0) circle (0.7);
    \node at (0,-1.0) {oogonium};
    \draw[-Stealth] (0.9,0) -- (1.5,0);
    \draw[omMeth, thick, fill=omMeth!20] (2.5,0.15) circle (0.65);
    \draw[omExo, thick, fill=omExo!20] (2.5,-0.7) circle (0.22);
    \node[font=\footnotesize] at (2.5,1.05) {egg};
    \node[font=\footnotesize] at (3.4,-0.7) {polar};
    \node[font=\footnotesize] at (3.4,-0.95) {body};
  \end{scope}
\end{tikzpicture}

{\small \omterm{def:g11:meiosis-life-cycles:gametogenesis}{Spermatogenesis} yields four small sperm; \omterm{def:g12:reproduction-development:gametogenesis}{oogenesis} typically
yields one large egg and \omterm{def:g11:meiosis-life-cycles:gametogenesis}{polar bodies} that discard extra \omterm{def:g11:meiosis-life-cycles:ploidy}{haploid} nuclei.}
\end{omfigure}

\begin{remark}[Timing in mammals]\label{rem:g12:reproduction-development:timing}
In many mammals, primary oocytes begin \omterm{def:g11:meiosis-life-cycles:meiosis}{meiosis} before birth and arrest;
cycles later complete maturation. Sperm are produced continuously after
puberty in human males.
\end{remark}

\section{Fertilisation}

\begin{definition}[Fertilisation]
\label{def:g12:reproduction-development:fertilisation}
\emph{Fertilisation}\index{fertilisation} is the fusion of sperm and egg
nuclei (and usually their cytoplasms) to form a \omterm{def:g12:reproduction-development:sexual}{zygote}. It restores the
\omterm{def:g11:meiosis-life-cycles:ploidy}{diploid} \omterm{def:g11:dna-replication:chromosome}{chromosome} number and activates \omterm{def:g10:scientific-biology:properties}{development}.
\end{definition}

\begin{method}[Steps of animal fertilisation]
\label{met:g12:reproduction-development:fert}
\begin{enumerate}
  \item Sperm contacts egg coats; species-recognition molecules act.
  \item Acrosome \omterm{def:g10:membranes-enzymes:enzyme}{enzymes} help penetrate outer layers.
  \item Membranes fuse; sperm nucleus enters.
  \item Blocks to polyspermy (fast electrical and/or slow cortical
        reaction) prevent multiple \omterm{def:g12:reproduction-development:fertilisation}{fertilisations}.
  \item \omterm{def:g11:meiosis-life-cycles:ploidy}{Haploid} nuclei unite as a \omterm{def:g11:meiosis-life-cycles:ploidy}{diploid} \omterm{def:g12:reproduction-development:sexual}{zygote} nucleus; \omterm{def:g12:reproduction-development:cleavage}{cleavage} begins.
\end{enumerate}
\end{method}

\begin{example}[Internal vs external fertilisation]
\label{ex:g12:reproduction-development:internal}
Many aquatic animals release \omterm{def:g12:reproduction-development:sexual}{gametes} into water (external). Terrestrial
vertebrates typically use internal \omterm{def:g12:reproduction-development:fertilisation}{fertilisation}, protecting \omterm{def:g12:reproduction-development:sexual}{gametes} from
desiccation and raising the chance of successful fusion.
\end{example}

\begin{definition}[External development and amniotic egg]
\label{def:g12:reproduction-development:amniote}
Reptiles, birds and monotremes often package the embryo in a shelled
\emph{amniotic egg}\index{amniotic egg} with membranes (amnion, chorion,
allantois, yolk sac) that manage water, \omterm{def:g12:animal-exchange:gas}{gas exchange} and waste --- a key
land \omterm{def:g11:evolution-mechanisms:fitness}{adaptation}. Placental mammals internalise analogous functions.
\end{definition}

\section{Cleavage and gastrulation}

\begin{definition}[Cleavage]\label{def:g12:reproduction-development:cleavage}
\emph{Cleavage}\index{cleavage} is the series of rapid early mitoses that
partition the \omterm{def:g12:reproduction-development:sexual}{zygote} into smaller \omterm{def:g10:the-cell:cell}{cells} (\emph{blastomeres}
\index{blastomere}) without much overall \omterm{def:g10:scientific-biology:properties}{growth}, producing a
\emph{blastula}\index{blastula} (or mammalian \omterm{def:g12:reproduction-development:human}{blastocyst}) with a fluid-filled
cavity.
\end{definition}

\begin{definition}[Gastrulation]
\label{def:g12:reproduction-development:gastrulation}
\emph{Gastrulation}\index{gastrulation} rearranges \omterm{def:g12:reproduction-development:cleavage}{blastula} \omterm{def:g10:the-cell:cell}{cells} into a
multilayered \emph{gastrula}\index{gastrula} with two or three
\emph{germ layers}\index{germ layer}: \omterm{prop:g12:reproduction-development:layers}{ectoderm}, \omterm{prop:g12:reproduction-development:layers}{endoderm} and (in
triploblasts) \omterm{prop:g12:reproduction-development:layers}{mesoderm}. \omterm{def:g10:the-cell:cell}{Cell} movements establish the body plan.
\end{definition}

\begin{proposition}[Fates of the germ layers]
\label{prop:g12:reproduction-development:layers}
In vertebrates, roughly: \emph{ectoderm}\index{ectoderm} $\to$ epidermis
and nervous system; \emph{endoderm}\index{endoderm} $\to$ gut lining and
associated organs; \emph{mesoderm}\index{mesoderm} $\to$ muscle, skeleton,
blood, kidney and much of the reproductive system. Details vary, but the
map organises \omterm{def:g12:reproduction-development:organo}{organogenesis}.
\end{proposition}
\admitted

\begin{omfigure}
\begin{tikzpicture}[font=\small]
  \draw[omDef, thick, fill=omDef!12] (0,0) circle (0.7);
  \node at (0,-1.15) {zygote};
  \draw[-Stealth] (1.0,0) -- (1.6,0);
  \draw[omDef, thick] (2.3,0.25) rectangle (3.1,0.85);
  \draw[omDef, thick] (2.3,-0.85) rectangle (3.1,-0.25);
  \draw[omDef, thick] (3.1,0.25) rectangle (3.9,0.85);
  \draw[omDef, thick] (3.1,-0.85) rectangle (3.9,-0.25);
  \node at (3.1,-1.15) {cleavage};
  \draw[-Stealth] (4.2,0) -- (4.8,0);
  \draw[omProp, thick, fill=omProp!15] (5.8,0) circle (0.75);
  \draw[omProp, thick, fill=white] (5.8,0) circle (0.35);
  \node at (5.8,-1.15) {blastula};
  \draw[-Stealth] (6.8,0) -- (7.4,0);
  \draw[omMeth, thick, fill=omMeth!20] (8.6,0) circle (0.75);
  \draw[omThm, thick, fill=omThm!20] (8.6,0) circle (0.45);
  \draw[omDef, thick, fill=omDef!25] (8.6,0) circle (0.2);
  \node at (8.6,-1.15) {gastrula};
\end{tikzpicture}

{\small Early \omterm{def:g10:scientific-biology:properties}{development}: \omterm{def:g12:reproduction-development:sexual}{zygote} $\to$ \omterm{def:g12:reproduction-development:cleavage}{cleavage} $\to$ \omterm{def:g12:reproduction-development:cleavage}{blastula} $\to$
\omterm{def:g12:reproduction-development:gastrulation}{gastrula} with layered organisation.}
\end{omfigure}

\begin{definition}[Organogenesis and neurulation]
\label{def:g12:reproduction-development:organo}
\emph{Organogenesis}\index{organogenesis} builds organs from \omterm{def:g12:reproduction-development:gastrulation}{germ layers}.
In chordates, \emph{neurulation}\index{neurulation} folds \omterm{prop:g12:reproduction-development:layers}{ectoderm} into
the neural tube, the rudiment of the CNS.
\end{definition}

\begin{remark}[Patterning]\label{rem:g12:reproduction-development:pattern}
Maternal determinants, inductive signals and gene-regulatory networks
assign \omterm{def:g10:the-cell:cell}{cell} \omterm{def:g10:molecules-of-life:lipids}{fates} along axes (anterior--posterior, dorsal--ventral).
Details occupy developmental \omterm{def:g10:scientific-biology:science}{biology}; the grammar is already visible at
\omterm{def:g12:reproduction-development:gastrulation}{gastrulation}.
\end{remark}

\section{Placental mammals}

\begin{definition}[Placenta]\label{def:g12:reproduction-development:placenta}
The \emph{placenta}\index{placenta} is an organ formed from embryonic
and maternal tissues that exchanges nutrients, gases and wastes between
maternal and fetal blood without free mixing of the two circulations.
It also secretes \omterm{def:g12:nerves-hormones:hormone}{hormones} that maintain pregnancy.
\end{definition}

\begin{definition}[Human early stages (outline)]
\label{def:g12:reproduction-development:human}
After \omterm{def:g12:reproduction-development:fertilisation}{fertilisation} in the oviduct, \omterm{def:g12:reproduction-development:cleavage}{cleavage} produces a
\emph{blastocyst}\index{blastocyst} that \emph{implants}\index{implantation}
in the uterine lining. The inner \omterm{def:g10:the-cell:cell}{cell} mass becomes the embryo; trophoblast
contributes to the \omterm{def:g12:reproduction-development:placenta}{placenta}. Gestation continues with \omterm{def:g10:scientific-biology:properties}{growth} and
organ maturation until birth.
\end{definition}

\begin{omfigure}
\begin{tikzpicture}[font=\small]
  \draw[omDef, thick, fill=omDef!10] (0,0) ellipse (2.4 and 1.3);
  \draw[omProp, thick, fill=omProp!20] (0.3,0.1) ellipse (1.1 and 0.7);
  \draw[omThm, thick, fill=omThm!25] (0.3,0.1) circle (0.35);
  \node[font=\footnotesize] at (0.3,0.1) {fetus};
  \draw[omMeth, very thick] (1.4,0.3) -- (2.8,1.0);
  \node[font=\footnotesize] at (2.5,1.35) {umbilical};
  \node[font=\footnotesize] at (-1.5,0) {uterus};
  \node[font=\footnotesize, omProp] at (0.3,-0.95) {placenta zone};
\end{tikzpicture}

{\small Schematic: fetus in the uterus exchanges with the mother via
\omterm{def:g12:reproduction-development:placenta}{placenta} and umbilical cord.}
\end{omfigure}

\begin{proposition}[Advantages of placentation]
\label{prop:g12:reproduction-development:placenta-adv}
Internal gestation with a \omterm{def:g12:reproduction-development:placenta}{placenta} protects the embryo, stabilises its
environment and allows extensive prenatal \omterm{def:g10:scientific-biology:properties}{growth} and brain \omterm{def:g10:scientific-biology:properties}{development}
in many mammals --- at energetic cost to the mother and with obstetric
constraints at birth.
\end{proposition}
\admitted

\begin{example}[Hormonal maintenance]
\label{ex:g12:reproduction-development:hcg}
In humans, early embryonic signals (e.g.\ hCG) maintain the corpus
luteum so \omterm{def:g12:reproduction-development:menstrual}{progesterone} keeps the uterine lining receptive until the
\omterm{def:g12:reproduction-development:placenta}{placenta} takes over \omterm{def:g12:nerves-hormones:hormone}{hormone} production.
\end{example}

\section{Menstrual cycle (outline)}

\begin{definition}[Ovarian and uterine cycles]
\label{def:g12:reproduction-development:menstrual}
In humans, the \emph{menstrual cycle}\index{menstrual cycle} coordinates
ovarian follicle maturation with preparation of the uterine lining.
\emph{FSH}\index{FSH} and \emph{LH}\index{LH} from the \omterm{def:g12:nerves-hormones:glands}{pituitary} drive
follicle \omterm{def:g10:scientific-biology:properties}{growth} and ovulation; ovarian \emph{oestrogen}\index{oestrogen}
and \emph{progesterone}\index{progesterone} rebuild and then maintain the
endometrium. If no pregnancy occurs, \omterm{def:g12:nerves-hormones:hormone}{hormone} withdrawal leads to
menstruation and the cycle restarts.
\end{definition}

\begin{omfigure}
\begin{tikzpicture}[font=\footnotesize]
  \draw[thick, -Stealth] (0,0) -- (8.2,0) node[right] {day};
  \draw[thick, -Stealth] (0,0) -- (0,3.0) node[above] {relative level};
  % FSH/LH sketch
  \draw[omDef, very thick, smooth]
    plot coordinates {(0.3,0.6) (1.5,0.9) (3.0,0.7) (3.8,2.4)
      (4.2,0.8) (6.5,0.5) (7.5,0.5)};
  \node[omDef] at (4.0,2.7) {LH surge};
  % oestrogen
  \draw[omMeth, very thick, smooth]
    plot coordinates {(0.3,0.4) (2.0,1.0) (3.5,1.8) (4.0,1.0)
      (5.5,1.4) (7.0,0.5)};
  \node[omMeth] at (2.2,1.5) {oestrogen};
  % progesterone
  \draw[omProp, very thick, smooth]
    plot coordinates {(0.3,0.25) (3.5,0.3) (4.5,0.4) (5.5,1.6)
      (6.5,1.5) (7.5,0.35)};
  \node[omProp] at (6.0,2.0) {progesterone};
  \draw[omExo, dashed] (3.9,0) -- (3.9,2.8);
  \node at (3.9,-0.35) {ovulation};
  \node at (1.8,-0.35) {follicular};
  \node at (6.0,-0.35) {luteal};
\end{tikzpicture}

{\small Schematic \omterm{def:g12:nerves-hormones:hormone}{hormone} sketch (not \omterm{def:g11:human-genetics:polygenic}{quantitative}): follicular \omterm{def:g12:reproduction-development:menstrual}{oestrogen}
rise, mid-cycle \omterm{def:g12:reproduction-development:menstrual}{LH} surge triggering ovulation, luteal \omterm{def:g12:reproduction-development:menstrual}{progesterone} high if
the corpus luteum is active.}
\end{omfigure}

\begin{proposition}[If fertilisation occurs]
\label{prop:g12:reproduction-development:pregnancy-signal}
\omterm{def:g12:reproduction-development:human}{Implantation} and early embryonic signals (e.g.\ hCG) maintain \omterm{def:g12:reproduction-development:menstrual}{progesterone}
output so the endometrium is not shed. Without that signal the corpus luteum
fails, \omterm{def:g12:reproduction-development:menstrual}{progesterone} falls, and menstruation begins --- resetting the uterus
for a possible later pregnancy.
\end{proposition}
\admitted

\begin{omfigure}
\begin{tikzpicture}[font=\footnotesize,
  box/.style={draw, rounded corners, minimum width=1.6cm,
    minimum height=0.65cm, align=center, fill=omDef!12}]
  \node[box] (z) at (0,0) {zygote};
  \node[box, fill=omMeth!15] (c) at (2.3,0) {cleavage};
  \node[box, fill=omProp!15] (b) at (4.6,0) {blastocyst};
  \node[box, fill=omThm!15] (i) at (6.9,0) {implant};
  \node[box, fill=omExo!15] (g) at (3.45,-1.4) {gastrula};
  \draw[-Stealth, thick] (z) -- (c);
  \draw[-Stealth, thick] (c) -- (b);
  \draw[-Stealth, thick] (b) -- (i);
  \draw[-Stealth, thick] (b.south) -- (g.north);
  \node[align=center] at (3.45,-2.2)
    {extra development timeline: layers form after implantation begins};
\end{tikzpicture}

{\small Human early stages in \omterm{def:g10:molecules-of-life:proteins}{sequence}: \omterm{def:g12:reproduction-development:fertilisation}{fertilisation} $\to$ \omterm{def:g12:reproduction-development:cleavage}{cleavage} $\to$
\omterm{def:g12:reproduction-development:human}{blastocyst} $\to$ \omterm{def:g12:reproduction-development:human}{implantation}, with \omterm{def:g12:reproduction-development:gastrulation}{gastrulation} establishing \omterm{def:g12:reproduction-development:gastrulation}{germ layers}.}
\end{omfigure}

\begin{method}[Comparing reproductive strategies]
\label{met:g12:reproduction-development:compare}
For any animal, ask: external or internal \omterm{def:g12:reproduction-development:fertilisation}{fertilisation}? Egg-laying or
live-bearing? How much yolk vs parental provisioning? How many offspring
and how much parental care? Trade-offs link to survival probability.
\end{method}

\section{Exercises}

\begin{exercise}[$\star$]\label{exo:g12:reproduction-development:1}
Define \omterm{def:g12:reproduction-development:sexual}{gamete}, \omterm{def:g12:reproduction-development:sexual}{zygote} and \omterm{def:g12:reproduction-development:fertilisation}{fertilisation}.
\end{exercise}

\begin{exercise}[$\star$]\label{exo:g12:reproduction-development:2}
How does the typical number of functional \omterm{def:g12:reproduction-development:sexual}{gametes} from one \omterm{def:g11:meiosis-life-cycles:meiosis}{meiosis}
differ between \omterm{def:g12:reproduction-development:gametogenesis}{spermatogenesis} and \omterm{def:g12:reproduction-development:gametogenesis}{oogenesis}?
\end{exercise}

\begin{exercise}[$\star$]\label{exo:g12:reproduction-development:3}
What is \omterm{def:g12:reproduction-development:cleavage}{cleavage}, and why does the embryo not grow much in mass during
early \omterm{def:g12:reproduction-development:cleavage}{cleavage}?
\end{exercise}

\begin{exercise}[$\star$]\label{exo:g12:reproduction-development:4}
Name the three \omterm{def:g12:reproduction-development:gastrulation}{germ layers} of a triploblastic animal.
\end{exercise}

\begin{exercise}[$\star\star$]\label{exo:g12:reproduction-development:5}
List the main steps that prevent polyspermy after the first sperm fuses.
\end{exercise}

\begin{exercise}[$\star\star$]\label{exo:g12:reproduction-development:6}
Assign nervous system, gut lining and skeletal muscle to their \omterm{def:g12:reproduction-development:gastrulation}{germ
layers}.
\end{exercise}

\begin{exercise}[$\star\star$]\label{exo:g12:reproduction-development:7}
What is \omterm{def:g12:reproduction-development:gastrulation}{gastrulation}, and why is it a critical transition in \omterm{def:g10:scientific-biology:properties}{development}?
\end{exercise}

\begin{exercise}[$\star\star$]\label{exo:g12:reproduction-development:8}
Explain how the \omterm{def:g12:reproduction-development:placenta}{placenta} exchanges substances without mixing maternal
and fetal blood.
\end{exercise}

\begin{exercise}[$\star\star$]\label{exo:g12:reproduction-development:9}
Why was the \omterm{def:g12:reproduction-development:amniote}{amniotic egg} a major innovation for life on land?
\end{exercise}

\begin{exercise}[$\star\star\star$]\label{exo:g12:reproduction-development:10}
Compare a frog (external \omterm{def:g12:reproduction-development:fertilisation}{fertilisation}, aquatic larva) with a human
(internal \omterm{def:g12:reproduction-development:fertilisation}{fertilisation}, \omterm{def:g12:reproduction-development:placenta}{placenta}) using the method of reproductive
strategies.
\end{exercise}

\begin{exercise}[$\star\star\star$]\label{exo:g12:reproduction-development:11}
A toxin blocks mitotic spindles during \omterm{def:g12:reproduction-development:cleavage}{cleavage}. Predict the consequence
for \omterm{def:g12:reproduction-development:human}{blastocyst} formation and \omterm{def:g12:reproduction-development:human}{implantation}.
\end{exercise}

\begin{exercise}[$\star\star\star$]\label{exo:g12:reproduction-development:12}
Why does \omterm{def:g11:meiosis-life-cycles:meiosis}{meiosis} matter for \omterm{def:g11:dna-replication:chromosome}{chromosome} number across generations, and
how do \omterm{prop:g11:meiosis-life-cycles:prophaseI}{crossing over} and independent assortment affect siblings?
\end{exercise}

\begin{exercise}[$\star\star$]\label{exo:g12:reproduction-development:13}
Describe the roles of \omterm{def:g12:reproduction-development:menstrual}{FSH}, \omterm{def:g12:reproduction-development:menstrual}{LH}, \omterm{def:g12:reproduction-development:menstrual}{oestrogen} and \omterm{def:g12:reproduction-development:menstrual}{progesterone} in a simplified
human \omterm{def:g12:reproduction-development:menstrual}{menstrual cycle}, including which event the \omterm{def:g12:reproduction-development:menstrual}{LH} surge triggers.
\end{exercise}

\begin{exercise}[$\star\star\star$]\label{exo:g12:reproduction-development:14}
Why does menstruation fail to occur if early pregnancy is established?
Link hCG (or equivalent signalling), the corpus luteum and the uterine
lining.
\end{exercise}

\section{Problem: From gametes to newborn}

\begin{problem}[{Weekend problem --- reconstruct the path from meiosis
to a placental newborn, with checkpoints for failure}]
\label{pb:g12:reproduction-development:1}
Trace human \omterm{def:g10:scientific-biology:properties}{reproduction} and early \omterm{def:g10:scientific-biology:properties}{development} as a \omterm{def:g10:molecules-of-life:proteins}{sequence} of
biological decisions.

\medskip
\noindent\textbf{Part I --- \omterm{def:g12:reproduction-development:sexual}{Gametes}.}
\begin{enumerate}
  \item State where sperm and eggs are produced and their ploidy.
  \item Why is unequal \omterm{def:g11:cell-cycle-mitosis:cytokinesis}{cytokinesis} in \omterm{def:g12:reproduction-development:gametogenesis}{oogenesis} adaptive for the egg?
  \item What genetic processes in \omterm{def:g11:meiosis-life-cycles:meiosis}{meiosis} make each \omterm{def:g12:reproduction-development:sexual}{gamete} unique?
\end{enumerate}

\medskip
\noindent\textbf{Part II --- \omterm{def:g12:reproduction-development:fertilisation}{Fertilisation} and early embryo.}
\begin{enumerate}[resume]
  \item Outline \omterm{def:g12:reproduction-development:fertilisation}{fertilisation} and the blocks to polyspermy.
  \item Describe \omterm{def:g12:reproduction-development:cleavage}{cleavage} up to the \omterm{def:g12:reproduction-development:human}{blastocyst} and name the two \omterm{def:g10:the-cell:cell}{cell}
        groups (inner \omterm{def:g10:the-cell:cell}{cell} mass vs trophoblast) and their \omterm{def:g10:molecules-of-life:lipids}{fates}.
  \item Define \omterm{def:g12:reproduction-development:human}{implantation}.
  \item What would multiple sperm nuclei joining one egg risk for
        \omterm{def:g11:dna-replication:chromosome}{chromosome} number?
\end{enumerate}

\medskip
\noindent\textbf{Part III --- \omterm{def:g12:reproduction-development:gastrulation}{Gastrulation} and layers.}
\begin{enumerate}[resume]
  \item What \omterm{def:g10:the-cell:cell}{cell} rearrangements achieve at \omterm{def:g12:reproduction-development:gastrulation}{gastrulation}?
  \item Give one adult derivative of each \omterm{def:g12:reproduction-development:gastrulation}{germ layer}.
  \item Link \omterm{def:g12:reproduction-development:organo}{neurulation} to the future brain and spinal cord.
\end{enumerate}

\medskip
\noindent\textbf{Part IV --- \omterm{def:g12:reproduction-development:placenta}{Placenta} and birth.}
\begin{enumerate}[resume]
  \item List three classes of substances that cross the \omterm{def:g12:reproduction-development:placenta}{placenta} and
        one that normally should not (in large amounts).
  \item Why is maternal health (nutrition, toxins, infection) critical
        during \omterm{def:g12:reproduction-development:organo}{organogenesis}?
  \item In one paragraph, contrast the \omterm{def:g12:reproduction-development:amniote}{amniotic egg} with the \omterm{def:g12:reproduction-development:placenta}{placenta}
        as solutions to embryonic life on land.
\end{enumerate}
\end{problem}
